% BEGIN PROFESSOR FEEDBACK PF20: Results - RQ2
\section{RQ2: Sensitivity to the Temporal Window}

AP and fingerprint coverage generally increase over the tested window range.
Figure~\ref{fig:res_rq2_delta} shows how held-out
average precision and motif coverage change across the nine tested windows in
the 40-role HuggingFace analysis. The baseline is fixed because its features do
not depend on $\delta$. Table~\ref{tab:res_rq2} shows the windows selected on the
development data in all three samples.

\begin{figure}[p]
  \centering
  \ResGraphic{0.48}{rq2_delta_performance_controlled_rep3}
  \ResGraphic{0.48}{rq2_delta_performance_prevalence_rep3}\\[4pt]
  \ResGraphic{0.48}{rq2_motif_coverage_controlled_rep3}
  \ResGraphic{0.48}{rq2_motif_coverage_prevalence_rep3}
  \caption{RQ2 temporal-window sensitivity in the 40-role HuggingFace analysis
  using the default timestamp ordering. The controlled cohort is shown on the
  left and the prevalence cohort on the right. The top row reports held-out
  average precision with shaded 95\,\% bootstrap confidence intervals.
  The bottom row reports fingerprint coverage across all addresses
  in each cohort.}
  \label{fig:res_rq2_delta}
\end{figure}

\IfFileExists{generated/results_rq2.tex}{\begin{table}[ht]
\centering
\caption{RQ2: development-CV window selections for both disjoint confirmations. The final column states whether the selected maximum was uniquely superior after corrected comparisons under all seven orderings.}
\label{tab:res_rq2}
\small
\begin{tabular}{@{}r l l l l@{}}
\toprule
\textbf{Rep.} & \textbf{Cohort} & \textbf{Maximum $\delta$} & \textbf{1-SE $\delta$} & \textbf{Unique maximum} \\
\midrule
1 & Controlled & 2 weeks & 1 week & no \\
1 & Prevalence & 30 days & 1 day & no \\
2 & Controlled & 2 weeks & 1 week & no \\
2 & Prevalence & 30 days & 30 days & no \\
\bottomrule
\end{tabular}
\end{table}
}{\ResPendingTable{RQ2 selected window against alternatives.}{tab:res_rq2}}

Both 32-role controlled cohorts select two weeks, while the 40-role controlled
cohort and all three prevalence cohorts select 30 days. In the 40-role analysis,
30 days is selected in both cohorts under all seven timestamp
orderings, while two weeks remains within one standard error of the maximum.
Mean development-CV AP increases from 0.7431 at five minutes to 0.7760 at 30
days in the controlled cohort, and from 0.5897 to 0.6253 in the prevalence
cohort. The binary Real-CATS comparison also selects 30 days, although three
days is already within one standard error of its maximum.

After correction for multiple comparisons, none of the three HuggingFace
samples has a selected maximum that is better than every other tested window
under all orderings. In the 40-role controlled cohort, five of the eight
comparisons with 30 days are significant using the default ordering. In the
prevalence cohort, none is significant. The 30-day window is the largest tested
value. It has the highest development score in the 40-role analysis, while
two weeks is within one standard error of that maximum.

Fingerprint coverage increases from 12.2\% at five minutes to 77.7\% at 30 days in
the 40-role controlled cohort. In the prevalence cohort, it increases
from 11.1\% to 75.9\%. This means that 22.3\% and 24.1\% of the respective addresses still
have no retained three-edge candidate pattern even at 30 days.
% END PROFESSOR FEEDBACK PF20

